\documentclass[aspectratio=169]{beamer}
\usetheme{Madrid}
\usecolortheme{default}
\usepackage[utf8]{inputenc}
\usepackage[T5]{fontenc}
\usepackage{graphicx}
\usepackage{hyperref}
\usepackage{booktabs}

\setbeamertemplate{caption}[numbered]
\renewcommand{\figurename}{Hình}

\title[Trí tuệ Nhân tạo cho Kế toán]{Trí tuệ Nhân tạo cho Kế toán \\ \vspace{0.3cm} \Large KẾ HOẠCH SLIDE LÝ THUYẾT - DAY 09 (BUỔI 9)}
\author{Đại học Đông Á}
\date{\today}

\begin{document}

% SLIDE 1
\begin{frame}
    \titlepage
\end{frame}

% SLIDE 2
\begin{frame}{Nội dung Bài học}
    \tableofcontents
\end{frame}


\section{KẾ HOẠCH SLIDE LÝ THUYẾT - DAY 09 (BUỔI 9)}


\begin{frame}{Title: Buổi 9: Ứng dụng AI trong Kiểm soát Nội bộ.}
    \begin{itemize}
        \item Phụ đề: Tự động hóa đánh giá KSNB và Khuôn khổ COSO 2013
        \item Giảng viên: Đại học Đông Á
        \item Môn học: Trí tuệ Nhân tạo cho Kế toán
    \end{itemize}
\end{frame}


\begin{frame}{Năng lực đạt được sau buổi học.}
    \begin{itemize}
        \item \textbf{Về Kiến thức:} Nắm vững 5 thành phần của Khuôn khổ COSO 2013. Hiểu được vai trò của AI trong việc tự động hóa đánh giá và thiết kế quy trình Kiểm soát Nội bộ (KSNB).
        \item \textbf{Về Kỹ năng:} Có khả năng sử dụng Prompt để yêu cầu AI rà soát lỗ hổng kiểm soát (như thiếu phân nhiệm, thiếu phê duyệt) trong một quy trình thực tế (ví dụ: quy trình mua hàng, thanh toán).
        \item \textbf{Về Tư duy:} Đánh giá tính rủi ro, phân biệt rõ vai trò không thể thay thế của con người (Professional Skepticism) so với sự hỗ trợ của AI trong môi trường KSNB.
    \end{itemize}
\end{frame}


\begin{frame}{Nội dung Chương trình (Agenda).}
    \begin{itemize}
        \item 1. Tổng quan về Kiểm soát Nội bộ và Khuôn khổ COSO 2013.
        \item 2. Ứng dụng AI trong Đánh giá và Nhận diện Rủi ro.
        \item 3. Tự động hóa Hoạt động Kiểm soát với AI.
        \item 4. Rà soát Lỗ hổng Quy trình (Thực tế: Quy trình Mua hàng).
        \item 5. Giới hạn của AI và Tương lai của KSNB.
        ### Phần 1: Tổng quan về KSNB và COSO 2013
    \end{itemize}
\end{frame}


\begin{frame}{Định nghĩa Kiểm soát Nội bộ (Internal Controls).}
    \begin{itemize}
        \item KSNB là một quy trình do Hội đồng quản trị, Ban giám đốc và các cá nhân khác thiết lập.
        \item Mục đích: Đảm bảo độ tin cậy của BCTC, tuân thủ pháp luật và hiệu quả hoạt động.
    \end{itemize}
\end{frame}


\begin{frame}{Khuôn khổ COSO 2013 là gì?}
    \begin{itemize}
        \item Committee of Sponsoring Organizations of the Treadway Commission (COSO).
        \item Là chuẩn mực toàn cầu được áp dụng rộng rãi nhất để thiết kế và đánh giá KSNB.
    \end{itemize}
\end{frame}


\begin{frame}{5 Thành phần của COSO.}
    \begin{itemize}
        \item 1. Môi trường kiểm soát (Control Environment).
        \item 2. Đánh giá rủi ro (Risk Assessment).
        \item 3. Hoạt động kiểm soát (Control Activities).
        \item 4. Thông tin và Truyền thông (Information \& Communication).
        \item 5. Giám sát (Monitoring).
    \end{itemize}
\end{frame}


\begin{frame}{Khối lập phương COSO (COSO Cube).}
    \begin{itemize}
        \item 3 Mục tiêu (Hoạt động, Báo cáo, Tuân thủ).
        \item 5 Thành phần KSNB.
        \item Cấu trúc tổ chức (Thực thể, Khối, Đơn vị, Chức năng).
    \end{itemize}
\end{frame}


\begin{frame}{Tại sao KSNB truyền thống lại gặp khó khăn?}
    \begin{itemize}
        \item Khối lượng giao dịch ngày càng khổng lồ.
        \item Rủi ro gian lận ngày càng tinh vi.
        \item Đánh giá KSNB thủ công tốn quá nhiều thời gian và chi phí.
    \end{itemize}
\end{frame}


\begin{frame}{AI - Giải pháp Tự động hóa.}
    \begin{itemize}
        \item Chuyển đổi từ "Kiểm tra chọn mẫu" sang "Kiểm tra toàn bộ 100\%".
        \item AI giúp phát hiện các khuôn mẫu (patterns) bất thường mà mắt người bỏ sót.
        ### Phần 2: Ứng dụng AI trong Đánh giá và Nhận diện Rủi ro
    \end{itemize}
\end{frame}


\begin{frame}{Nhận diện Rủi ro (Risk Identification).}
    \begin{itemize}
        \item Theo COSO, rủi ro là các sự kiện có tác động tiêu cực đến mục tiêu kinh doanh.
        \item AI sử dụng Xử lý ngôn ngữ tự nhiên (NLP) để quét các hợp đồng, email, báo cáo nhằm tìm rủi ro tiềm ẩn.
    \end{itemize}
\end{frame}


\begin{frame}{Phân tích Rủi ro Gian lận (Fraud Risk).}
    \begin{itemize}
        \item Gian lận BCTC, Biển thủ tài sản, Tham nhũng.
        \item AI có khả năng lập hồ sơ rủi ro dựa trên dữ liệu lịch sử và học máy (Machine Learning).
    \end{itemize}
\end{frame}


\begin{frame}{Nhận diện Dấu hiệu Cảnh báo (Red Flags).}
    \begin{itemize}
        \item Thay vì đợi cuối tháng/cuối năm mới kiểm tra.
        \item AI giúp cảnh báo theo thời gian thực (Real-time alerts).
    \end{itemize}
\end{frame}


\begin{frame}{Phân tích Rủi ro Người nội bộ (Insider Threat).}
    \begin{itemize}
        \item Phân tích hành vi người dùng (UBA - User Behavior Analytics).
        \item Ví dụ: Kế toán viên đăng nhập hệ thống ERP vào lúc 2h sáng chủ nhật.
    \end{itemize}
\end{frame}


\begin{frame}{Sự bất phân nhiệm (Segregation of Duties - SoD).}
    \begin{itemize}
        \item Nguyên tắc cốt lõi: Một người không được kiêm nhiệm cả việc tạo giao dịch, phê duyệt và hạch toán.
    \end{itemize}
\end{frame}


\begin{frame}{AI rà soát Phân nhiệm tự động.}
    \begin{itemize}
        \item AI quét hệ thống phân quyền của ERP để tìm ra những user có "quyền lực kép" gây nguy cơ gian lận.
    \end{itemize}
\end{frame}


\begin{frame}{Đánh giá Rủi ro Động (Dynamic Risk Assessment).}
    \begin{itemize}
        \item Rủi ro không tĩnh, nó thay đổi hàng ngày theo thị trường.
        \item AI tích hợp dữ liệu bên ngoài (tin tức, biến động giá) để cập nhật hồ sơ rủi ro.
        ### Phần 3: Tự động hóa Hoạt động Kiểm soát
    \end{itemize}
\end{frame}


\begin{frame}{Hoạt động Kiểm soát (Control Activities).}
    \begin{itemize}
        \item Là các chính sách và thủ tục để đảm bảo các chỉ thị giảm thiểu rủi ro được thực hiện.
        \item Gồm kiểm soát phòng ngừa (Preventive) và kiểm soát phát hiện (Detective).
    \end{itemize}
\end{frame}


\begin{frame}{Tự động hóa Kiểm soát Phòng ngừa.}
    \begin{itemize}
        \item AI chặn ngay lập tức một hóa đơn thanh toán nếu số tài khoản nhà cung cấp mới được thay đổi trong vòng 24h.
    \end{itemize}
\end{frame}


\begin{frame}{Tự động hóa Kiểm soát Phát hiện.}
    \begin{itemize}
        \item AI tự động đối chiếu sổ phụ ngân hàng (Bank Reconciliation) với hàng vạn dòng giao dịch trong vài giây.
    \end{itemize}
\end{frame}


\begin{frame}{Kế toán Không mã code (No-code Accounting).}
    \begin{itemize}
        \item Sử dụng ChatGPT/Copilot để tạo ra các quy trình tự động hóa mà không cần biết lập trình (Python/SQL).
    \end{itemize}
\end{frame}


\begin{frame}{Rà soát Hợp đồng bằng NLP.}
    \begin{itemize}
        \item Kiểm tra xem hợp đồng mua bán có thiếu các điều khoản phạt vi phạm hay không.
    \end{itemize}
\end{frame}


\begin{frame}{Kiểm tra Hóa đơn tự động (OCR + AI).}
    \begin{itemize}
        \item OCR đọc hóa đơn, AI đối chiếu 3 chiều (Three-way matching) giữa: Đơn đặt hàng (PO) - Phiếu nhập kho - Hóa đơn.
    \end{itemize}
\end{frame}


\begin{frame}{Khám phá Lỗ hổng Kế toán.}
    \begin{itemize}
        \item Tự động quét sổ cái để tìm các nghiệp vụ Nợ/Có lặp lại với số tiền dưới mức phê duyệt (Structuring fraud).
    \end{itemize}
\end{frame}


\begin{frame}{Nâng cấp Thông tin và Truyền thông.}
    \begin{itemize}
        \item Chatbot nội bộ trả lời các câu hỏi về chính sách kế toán cho nhân viên.
        ### Phần 4: Rà soát Lỗ hổng Quy trình (Thực tế)
    \end{itemize}
\end{frame}


\begin{frame}{Tình huống: Quy trình Mua hàng.}
    \begin{itemize}
        \item Nhân viên mua hàng tự liên hệ nhà cung cấp, tự nhận hàng và tự làm đề nghị thanh toán gửi Kế toán trưởng.
    \end{itemize}
\end{frame}


\begin{frame}{Sử dụng AI để tìm Lỗ hổng.}
    \begin{itemize}
        \item Kế toán viên có thể "nạp" mô tả quy trình trên vào ChatGPT và yêu cầu tìm rủi ro.
    \end{itemize}
\end{frame}


\begin{frame}{Các lỗ hổng do AI phát hiện.}
    \begin{itemize}
        \item Thiếu phê duyệt từ Trưởng bộ phận trước khi mua.
        \item Thiếu phân nhiệm giữa việc Mua hàng và Nhận hàng.
    \end{itemize}
\end{frame}


\begin{frame}{Thiết kế lại Quy trình (Re-design).}
    \begin{itemize}
        \item Yêu cầu AI đề xuất một quy trình mua hàng mới, chuẩn COSO.
    \end{itemize}
\end{frame}


\begin{frame}{Lưu đồ Kiểm soát Mới.}
    \begin{itemize}
        \item Tách biệt: Người đề xuất - Người phê duyệt - Người mua - Người nhận hàng - Kế toán thanh toán.
    \end{itemize}
\end{frame}


\begin{frame}{Đưa AI vào Giám sát.}
    \begin{itemize}
        \item AI đóng vai trò là "Người giám sát ảo" liên tục audit quá trình này.
    \end{itemize}
\end{frame}


\begin{frame}{Tạo Prompt Kiểm toán Tự động.}
    \begin{itemize}
        \item Prompt: "Đóng vai chuyên gia COSO, hãy tìm 3 rủi ro lớn nhất trong quy trình thanh toán sau đây..."
    \end{itemize}
\end{frame}


\begin{frame}{Tối ưu hóa Câu hỏi Prompt.}
    \begin{itemize}
        \item Kế toán phải cung cấp đủ context: Loại hình công ty, quy mô, quy chế chi tiêu nội bộ.
        ### Phần 5: Giới hạn của AI và Tương lai của KSNB
    \end{itemize}
\end{frame}


\begin{frame}{Không có AI nào hoàn hảo.}
    \begin{itemize}
        \item Hiện tượng "Ảo giác" (Hallucination) của AI: AI bịa ra các lỗi không có thật trong hệ thống.
    \end{itemize}
\end{frame}


\begin{frame}{Sự Cố hữu của Rủi ro.}
    \begin{itemize}
        \item Ban giám đốc vẫn có thể "vượt mặt" kiểm soát (Management Override of Controls). AI rất khó phát hiện nếu sếp thông đồng.
    \end{itemize}
\end{frame}


\begin{frame}{Vấn đề Bảo mật Dữ liệu.}
    \begin{itemize}
        \item Khi đưa dữ liệu nội bộ lên ChatGPT, công ty đối mặt rủi ro lộ lọt thông tin.
    \end{itemize}
\end{frame}


\begin{frame}{Tính Thiên lệch (Bias).}
    \begin{itemize}
        \item AI có thể tạo ra các cảnh báo giả (False Positives) đối với các giao dịch hợp lệ.
    \end{itemize}
\end{frame}


\begin{frame}{Vai trò của Con người (Human in the loop).}
    \begin{itemize}
        \item AI chỉ đóng vai trò "Trợ lý", quyết định cuối cùng vẫn thuộc về Kiểm toán viên nội bộ và Kế toán trưởng.
    \end{itemize}
\end{frame}


\begin{frame}{Sự Hoài nghi Nghề nghiệp (Professional Skepticism).}
    \begin{itemize}
        \item Kế toán viên phải luôn đặt câu hỏi về độ tin cậy của báo cáo rủi ro do AI xuất ra.
    \end{itemize}
\end{frame}


\begin{frame}{Tương lai của KSNB.}
    \begin{itemize}
        \item Từ kiểm toán định kỳ (Periodic) sang kiểm toán liên tục (Continuous Auditing).
    \end{itemize}
\end{frame}


\begin{frame}{Kỹ năng mới của Kế toán viên.}
    \begin{itemize}
        \item Không chỉ hạch toán, mà còn là người thiết kế hệ thống KSNB tự động và "điều khiển" AI.
    \end{itemize}
\end{frame}


\begin{frame}{Tổng kết.}
    \begin{itemize}
        \item COSO 2013 + AI = Hệ thống kiểm soát vững chắc, nhanh chóng và thông minh.
    \end{itemize}
\end{frame}


\begin{frame}{Hỏi và Đáp (Q\&A).}
    \begin{itemize}
        \item Giải đáp thắc mắc về ứng dụng AI trong kiểm soát thực tế tại doanh nghiệp.
    \end{itemize}
\end{frame}


\end{document}